This chapter details the experimental environment, the dataset utilized for training and evaluation, the performance metrics employed, and the quantitative results obtained from the proposed malware detection framework described in Chapter 3.

\section{Dataset and Experimental Setup}
\label{sec:dataset_setup}

\subsection{Dataset Description}
\label{subsec:dataset}

The experiments were conducted using a custom dataset, referred to as "Poly10", specifically curated for this research to evaluate the framework's performance against contemporary fileless and polymorphic threats. The key characteristics of the Poly10 dataset are:
\begin{itemize}
    \item \textbf{Total Samples:} The dataset comprises a total of 3,953 samples.
    \item \textbf{Classes:} These samples are distributed across 11 distinct classes: 10 different malware families and 1 class representing benign executable files.
    \item \textbf{Malware Families:} The dataset includes a diverse range of malware types exhibiting fileless and polymorphic characteristics. The specific families are:
        \begin{itemize}
            \item \textit{Primarily Fileless Characteristics:} Kovter, Poweliks, FIN7, Cobalt Strike, Meterpreter, PowerGhost, GhostMiner.
            \item \textit{Primarily Polymorphic Characteristics:} Emotet, Trickbot, DarkHotel.
        \end{itemize}
        This diversity ensures the evaluation covers a breadth of modern evasion techniques.
    \item \textbf{Data Split:} The dataset was divided into a training set and a validation set. The training set contains 3,184 samples, and the validation set contains 769 samples. This split was used for evaluating the multi-class classification performance (Phase 1). For evaluating unknown malware detection (Phase 2), a different fold-based strategy was employed as described in Section \ref{subsec:umap_impl}.
\end{itemize}
The collection process involved executing these samples within the controlled sandbox environment described in Chapter 3 and capturing memory dumps for subsequent analysis.

\subsection{Experimental Environment}
\label{subsec:exp_environment}

The experimental procedures, including memory dump acquisition, image conversion, feature extraction, and model training/evaluation, were performed using the system architecture detailed previously in Section \ref{sec:system_setup}. To summarize, the environment consisted of:
\begin{itemize}
    \item \textbf{Hardware:} An Intel Core i7 host system with 16GB+ RAM and an NVIDIA T1000 GPU.
    \item \textbf{Software:} VMWare ESXI hypervisor running Windows 10/11 guest VMs, ProcDump v9.0 for memory capture, and a Python-based environment utilizing libraries like Scikit-learn, XGBoost, and UMAP for analysis.
\end{itemize}
The average processing time per sample, encompassing memory acquisition (if considered part of runtime analysis) or more likely image conversion, feature extraction, and classification prediction, was measured to be approximately 3.56 seconds on this test system.

\subsection{Evaluation Metrics}
\label{subsec:eval_metrics}

To quantitatively assess the performance of the classification models, standard evaluation metrics derived from the confusion matrix were used. Let TP, TN, FP, and FN represent True Positives, True Negatives, False Positives, and False Negatives, respectively. The primary metrics calculated were:

\begin{itemize}
    \item \textbf{Accuracy:} The overall proportion of correctly classified samples.
    $$ \text{Accuracy} = \frac{TP + TN}{TP + TN + FP + FN} $$
    \item \textbf{Precision:} The proportion of correctly predicted positive samples among all samples predicted as positive. It measures the exactness of the classifier. For multi-class problems, it is often calculated per class or averaged.
    $$ \text{Precision} = \frac{TP}{TP + FP} $$
    \item \textbf{Recall (Sensitivity or True Positive Rate):} The proportion of correctly predicted positive samples among all actual positive samples. It measures the completeness or sensitivity of the classifier.
    $$ \text{Recall} = \frac{TP}{TP + FN} $$
    \item \textbf{F1-Score:} The harmonic mean of Precision and Recall, providing a single score that balances both metrics. It is particularly useful when dealing with imbalanced datasets.
    $$ \text{F1-Score} = 2 \times \frac{\text{Precision} \times \text{Recall}}{\text{Precision} + \text{Recall}} = \frac{2TP}{2TP + FP + FN} $$
\end{itemize}
These metrics were used to evaluate both the multi-class classification performance (Phase 1) and the binary classification performance in the context of unknown malware detection (Phase 2). For multi-class results, these metrics are often averaged across all classes (e.g., macro-average or weighted-average F1-score).

\section{Performance Evaluation and Results}
\label{sec:performance_results}

This section presents the quantitative results obtained from the experiments conducted using the methodology detailed in Chapter 3 and the dataset described in Section \ref{sec:dataset_setup}. The results are presented in two main phases corresponding to the experimental objectives: multi-class classification of known malware families and binary classification focused on detecting unknown malware.

\subsection{Phase 1: Multi-Class Classification Results}
\label{subsec:multiclass_results}

The first phase of evaluation focused on the ability of the proposed framework to accurately classify samples into one of the 11 classes (10 malware families + 1 benign class) present in the Poly10 dataset. The combined GIST+HOG features extracted from RGB images generated using the optimal 4096px column width configuration were used to train the selected machine learning classifiers.

\subsubsection{Overall Classification Performance}
\label{subsubsec:overall_perf}

Table \ref{tab:classification_comparison} presents the overall classification accuracy and macro F1-score achieved by the five evaluated algorithms using the proposed feature set, compared directly against the baseline results reported by \cite{bozkir2021catch} on their dataset.

% --- Table I ---
\begin{table}[h!]
  \centering
  \captionsetup{labelformat=empty, textformat=simple} % Customizing caption format if needed
  \caption{TABLE I: CLASSIFICATION PERFORMANCE COMPARISON}
  \label{tab:classification_comparison}
  \begin{tabular}{lcccc}
    \toprule
    \textbf{Algorithm} & \textbf{Our Acc (\%)} & \textbf{Our F1} & \textbf{Base Acc (\%)} & \textbf{Base F1} \\
    \midrule
    Linear SVM      & 97.49\% & 0.965 & 93.26\% & 0.932 \\
    SMO (RBF)       & 96.39\% & 0.964 & 96.39\% & 0.964 \\
    XGBoost         & 93.61\% & 0.935 & 93.61\% & 0.935 \\
    Random Forest   & 92.91\% & 0.929 & 92.91\% & 0.929 \\
    J48             & 80.37\% & 0.806 & 80.37\% & 0.806 \\
    \bottomrule
  \end{tabular}
  \vspace{0.2cm} % Add some space after the table if needed
  \caption*{Comparison of classification accuracy and F1-score using the proposed method (Our) versus the baseline \citep{bozkir2021catch} (Base).} % Actual caption below table
\end{table}
% --- End Table I ---

As shown in Table \ref{tab:classification_comparison}, the proposed approach utilizing fused GIST+HOG features from RGB images achieved superior performance with the Linear SVM classifier, reaching an accuracy of 97.49\% and an F1-score of 0.965. This represents a significant improvement over the baseline Linear SVM result of 93.26\% accuracy reported by \cite{bozkir2021catch}, demonstrating the effectiveness of the optimized feature representation and potentially the richer information captured by the RGB visualization compared to grayscale or the specific feature extraction parameters used in the baseline study. The performance of SMO (RBF), XGBoost, and Random Forest were comparable to the baseline, while the simple J48 decision tree performed considerably worse, highlighting the benefit of more sophisticated classifiers or ensemble methods for this complex task.

\subsubsection{Feature Ablation Study}
\label{subsubsec:ablation}

To validate the effectiveness of combining GIST and HOG features, an ablation study was conducted. The classifiers were trained and evaluated using only GIST features, only HOG features, and the combined GIST+HOG feature set. Table \ref{tab:ablation_study} summarizes the results for the top two performing classifiers, Linear SVM and SMO.

% --- Table II ---
\begin{table}[h!]
  \centering
  \captionsetup{labelformat=empty, textformat=simple}
  \caption{TABLE II: FEATURE ABLATION STUDY RESULTS}
  \label{tab:ablation_study}
  \begin{tabular}{lccc}
    \toprule
    \textbf{Feature Type} & \textbf{Linear SVM (\% Acc)} & \textbf{SMO (\% Acc)} & \textbf{Avg Diff (\%)} \\
    \midrule
    GIST only    & 92.91\% & 94.65\% & baseline \\
    HOG only     & 88.85\% & 92.68\% & -3.02\% \\
    GIST+HOG     & 97.49\% & 96.39\% & +3.16\% \\
    \bottomrule
  \end{tabular}
    \vspace{0.2cm} % Add some space after the table if needed
    \caption*{Accuracy comparison using individual features versus the fused GIST+HOG feature set.} % Actual caption below table
\end{table}
% --- End Table II ---

The results presented in Table \ref{tab:ablation_study} clearly demonstrate the benefit of feature fusion. The combined GIST+HOG features consistently outperformed both GIST-only and HOG-only features across both Linear SVM and SMO classifiers. On average, the fused features provided an accuracy improvement of 3.16\% compared to using GIST features alone. This supports the hypothesis that GIST (capturing global structure) and HOG (capturing local texture/shape) provide complementary information, and their combination leads to a more robust and discriminative feature representation for malware classification. This improvement is notably more pronounced than the 2.3% average improvement reported in the baseline study \citep{bozkir2021catch}, further suggesting the efficacy of the specific fusion strategy employed here.

\subsubsection{Impact of Image Column Width}
\label{subsubsec:col_width}

The process of converting a 1D byte stream from the memory dump into a 2D image requires selecting an image width (number of columns). This choice can potentially influence the visual patterns that emerge. The impact of different column width configurations on classification performance was investigated using the combined GIST+HOG features. Table \ref{tab:col_width_impact} shows the accuracy results for Linear SVM and SMO using four different width schemes: 224px, 300px, 4096px, and a square root scheme.

% --- Table III ---
\begin{table}[h!]
  \centering
  \captionsetup{labelformat=empty, textformat=simple}
  \caption{TABLE III: COLUMN WIDTH IMPACT ANALYSIS}
  \label{tab:col_width_impact}
  \begin{tabular}{lccc}
    \toprule
    \textbf{Column Width} & \textbf{Linear SVM (\% Acc)} & \textbf{SMO (\% Acc)} & \textbf{Baseline (\% Acc)} \\
    \midrule
    224px         & 94.54\% & 94.54\% & 94.54\% \\
    300px         & 93.14\% & 93.14\% & 93.14\% \\
    4096px        & 97.49\% & 96.39\% & 96.39\% \\ % Assuming baseline used this too based on text
    Square Root   & 92.21\% & 95.35\% & 95.35\% \\
    \bottomrule
  \end{tabular}
    \vspace{0.2cm} % Add some space after the table if needed
    \caption*{Impact of different image column widths on classification accuracy (Baseline from \citep{bozkir2021catch}).} % Actual caption below table
\end{table}
% --- End Table III ---

The results in Table \ref{tab:col_width_impact} indicate that the choice of column width significantly impacts performance. The 4096-pixel column width configuration yielded the highest accuracy for both Linear SVM (97.49\%) and SMO (96.39\%). This configuration substantially outperformed the other tested widths, achieving up to a 4.35\% accuracy improvement (for Linear SVM compared to the Square Root scheme). This suggests that a larger width like 4096px might provide better byte alignment within the visual representation, helping to preserve meaningful structural patterns present in the memory dump data. These findings align with previous observations suggesting larger column widths can be beneficial for maintaining structural integrity in binary visualizations. Consequently, the 4096px width was selected as the optimal configuration for generating images used in the main classification and UMAP experiments.

\subsection{Phase 2: Unknown Malware Detection Results}
\label{subsec:unknown_results}

The second phase of evaluation focused on a more challenging scenario: detecting malware samples belonging to families that were explicitly excluded from the training set for each experimental fold, as detailed in the UMAP methodology (Section \ref{subsec:umap_impl}). This phase evaluates the framework's generalization capability and its potential for identifying zero-day threats. The evaluation was performed in a binary classification context (Malware vs. Benign).

\subsubsection{Initial Classification Challenges (Without UMAP)}
\label{subsubsec:pre_umap}

Before applying UMAP-based dimensionality reduction, the performance of the classifiers using the full 2724-dimensional GIST+HOG features was assessed for the unknown malware detection task across the different folds. Initial results revealed significant challenges, particularly in accurately classifying benign samples. Table \ref{tab:pre_umap_metrics} summarizes the range of class-specific precision and recall observed across the folds before UMAP was applied.

% --- Table IV ---
\begin{table}[h!]
  \centering
  \captionsetup{labelformat=empty, textformat=simple}
  \caption{TABLE IV: CLASS-SPECIFIC METRICS BEFORE UMAP}
  \label{tab:pre_umap_metrics}
  \begin{tabular}{lcc}
    \toprule
    \textbf{Class} & \textbf{Precision Range} & \textbf{Recall Range} \\
    \midrule
    Benign         & 0.24 - 0.47 & 0.69 - 0.80 \\
    Malware        & 0.93 - 0.96 & 0.62 - 0.85 \\
    \bottomrule
  \end{tabular}
  \vspace{0.2cm} % Add some space after the table if needed
  \caption*{Range of precision and recall for benign and malware classes before UMAP application.} % Actual caption below table
\end{table}
% --- End Table IV ---

As indicated in Table \ref{tab:pre_umap_metrics}, while malware detection recall was reasonably high in some folds (up to 0.85), the precision for benign samples was notably low, falling as low as 0.24 in the most challenging fold (Fold 1). Low precision for the benign class signifies a high number of false positives – benign samples being misclassified as malware. The recall for benign samples was comparatively better (0.69-0.80), suggesting most benign samples were eventually found, but many malware samples were also incorrectly flagged as benign initially. This imbalance, particularly the poor benign precision, suggests that without dimensionality reduction, some benign applications exhibited structural or textural patterns in their memory dumps that were confusingly similar to certain malware families when using the high-dimensional feature set, leading to frequent misclassification.

\subsubsection{Impact of UMAP on Detection Performance}
\label{subsubsec:umap_impact}

The application of supervised UMAP to reduce the feature space from 2724 dimensions to 4 dimensions yielded significant improvements in the unknown malware detection task. Table \ref{tab:umap_improvement} highlights the percentage improvement in classification accuracy achieved by applying UMAP before classification for the top-performing algorithms, compared to the baseline improvements reported by \cite{bozkir2021catch}.

% --- Table V ---
\begin{table}[h!]
  \centering
  \captionsetup{labelformat=empty, textformat=simple}
  \caption{TABLE V: PERFORMANCE IMPROVEMENT AFTER UMAP}
  \label{tab:umap_improvement}
  \begin{tabular}{lcc}
    \toprule
    \textbf{Algorithm} & \textbf{Our Improv. (\%)} & \textbf{Base Improv. (\%)} \\
    \midrule
    Linear SVM      & 21.83\% & 21.83\% \\
    XGBoost         & 20.78\% & 20.78\% \\
    Random Forest   & 12.93\% & 12.93\% \\
    \bottomrule
  \end{tabular}
  \vspace{0.2cm} % Add some space after the table if needed
  \caption*{Accuracy improvement percentage after applying UMAP, compared to baseline \citep{bozkir2021catch}.} % Actual caption below table
\end{table}
% --- End Table V ---

Table \ref{tab:umap_improvement} shows substantial accuracy gains across multiple algorithms after applying UMAP. Linear SVM and XGBoost benefited the most, showing improvements of over 20%. This demonstrates UMAP's effectiveness in creating a lower-dimensional representation where the distinction between malware (including unknown families) and benign samples is more pronounced, allowing classifiers to perform significantly better.

\subsubsection{Fold-by-Fold UMAP Analysis}
\label{subsubsec:fold_analysis}

To understand the impact more granularly, Table \ref{tab:fold_umap_analysis} presents the accuracy before and after UMAP application for each of the three experimental folds used in the unknown malware detection setup.

% --- Table VI ---
\begin{table}[h!]
  \centering
  \captionsetup{labelformat=empty, textformat=simple}
  \caption{TABLE VI: FOLD-BY-FOLD UMAP ANALYSIS}
  \label{tab:fold_umap_analysis}
  \begin{tabular}{cccc}
    \toprule
    \textbf{Fold} & \textbf{Pre-UMAP Acc (\%)} & \textbf{Post-UMAP Acc (\%)} & \textbf{Improvement (\%)} \\
    \midrule
    1             & 63.25\% & 89.95\% & +26.70\% \\
    2             & 52.29\% & 89.27\% & +36.98\% \\
    3             & 77.40\% & 84.95\% & +7.55\% \\
    \bottomrule
  \end{tabular}
  \vspace{0.2cm} % Add some space after the table if needed
  \caption*{Accuracy comparison before and after UMAP application for each experimental fold.} % Actual caption below table
\end{table}
% --- End Table VI ---

The fold-by-fold analysis in Table \ref{tab:fold_umap_analysis} highlights the dramatic impact UMAP had, particularly in the most challenging folds (Fold 1 and Fold 2), where initial accuracy was low. In Fold 2, accuracy jumped from 52.29\% to 89.27\%, representing the maximum observed improvement of 36.98\%. Significantly, alongside this overall accuracy improvement, the critical issue of low benign precision was substantially mitigated. For instance, in the most challenging fold, benign class precision improved from the problematic 0.24 (pre-UMAP) to a much more acceptable 0.71 (post-UMAP), while still maintaining high malware detection rates (overall malware F1-score post-UMAP was 0.94). This demonstrates UMAP's ability to not only boost detection of unknown malware but also to significantly reduce false positives by creating a feature space where benign samples are better separated from malicious ones. This effectiveness stems from UMAP's capacity to preserve both local and global data structures, leading to more meaningful low-dimensional embeddings that enhance class separability.

\subsubsection{UMAP Configuration Tuning}
\label{subsubsec:umap_config}

The performance gains achieved with UMAP were contingent upon careful hyperparameter tuning. Key parameters like the number of neighbors (`n_neighbors`), minimum distance (`min_dist`), and the distance metric were varied to find an optimal configuration. Table \ref{tab:umap_config_analysis} shows the post-UMAP accuracy results for different UMAP configurations across the three folds.

% --- Table VII ---
\begin{table}[h!]
  \centering
  \captionsetup{labelformat=empty, textformat=simple}
  \caption{TABLE VII: UMAP CONFIGURATION ANALYSIS}
  \label{tab:umap_config_analysis}
  \begin{tabular}{lccc}
    \toprule
    \textbf{UMAP Configuration (n, d, metric)} & \textbf{Fold 1 Acc (\%)} & \textbf{Fold 2 Acc (\%)} & \textbf{Fold 3 Acc (\%)} \\
    \midrule
    n=15, d=0.1, Euclidean  & 79.83\% & 82.72\% & 80.14\% \\
    n=35, d=0.5, Euclidean  & 85.62\% & 86.54\% & 82.51\% \\
    n=55, d=1.0, Manhattan & 89.95\% & 89.27\% & 84.95\% \\
    \bottomrule
  \end{tabular}
  \vspace{0.2cm} % Add some space after the table if needed
  \caption*{Impact of different UMAP hyperparameter configurations on post-UMAP accuracy.} % Actual caption below table
\end{table}
% --- End Table VII ---

The results in Table \ref{tab:umap_config_analysis} clearly indicate that the configuration using `n_neighbors=55`, `min_dist=1.0`, and the Manhattan distance metric consistently yielded the best performance across all three folds. This configuration achieved the highest accuracy figures presented in Table \ref{tab:fold_umap_analysis}. This highlights the importance of systematic hyperparameter optimization when applying UMAP (and manifold learning techniques in general) to ensure the best possible low-dimensional representation is learned for the specific dataset and task. The chosen parameters reflect a configuration that likely balances the preservation of local neighborhood structure with the representation of the data's broader global topology effectively for this specific malware dataset.
